\documentclass[11pt,a4paper]{article}
\usepackage[margin=22mm]{geometry}
\usepackage{amsmath,amssymb}
\usepackage[hidelinks]{hyperref}
\setlength{\parindent}{0pt}
\setlength{\parskip}{5pt}
\raggedright
\begin{document}
{\Large Aligning the update without losing retained traffic counts}\\[4pt]
8 October 2026 --- K.A approved; field/count interface F proposed

\textbf{Scope.} Conditional own-field bookkeeping, two analytical diagnostics,
and a constructive local reconstruction. Restriction records carry zero extra
traffic mass. No full endogenous solver, reached metering/TST comparison,
clearance theorem or gradient claim is inferred. Primary-source boundaries
are in \texttt{joined\_step\_source\_audit.md}.

\section*{1. The selected timing uses two different outputs}
Let an old record start at $x_j^k$ with sampled desired speed $s_j^k$.
Write $c_j(\tau)$ for equilibrium speed of its own one-sided ahead traffic
at elapsed time $\tau$. Conditional on that jointly determined trace, K.A is
\begin{equation}
 v_j(\tau)=\min\left\{s_j^k,\inf_{0\le r\le\tau}c_j(r)\right\},\quad
 x_j^{k+1}=x_j^k+\int_0^{\Delta t}v_j(\tau)\,d\tau,\quad
 v_{j,\mathrm{end}}=v_j(\Delta t-).
\end{equation}
The essential infimum can be used for merely measurable traces; one-sided
contact values determine braking. Positive recovery is refreshed only at the
next regular update, by the selected source recurrence from ending speed,
with its existing old-position grade convention. Interior release does not
increase $v_j$ in this interval. This formal phase rule is not an existence
proof for the own trace. E1 births follow completed transfers and start acting
in the next interval.

\section*{2. Priority belongs to the same receiving history}
For ramp and mainline demands $d_r,d_m$ and shared receiving rate $R$, all in
equivalents/s, fixed ramp priority $p\in(0,1)$ with reallocation gives
\begin{equation}
 f_r=\min\{d_r,\max(pR,R-d_m)\},\qquad
 f_m=\min\{d_m,\max((1-p)R,R-d_r)\}.
\end{equation}
Thus $0\le f_a\le d_a$ and $f_r+f_m=\min(d_r+d_m,R)$. If requests evolve,
accepted masses are
\begin{equation}
 T_a=\int_{t_k}^{t_{k+1}} f_a(d_r(t),d_m(t),R(t))\,dt.
\end{equation}
Allocation and time integration generally do not commute. With $p=1/2$,
two one-second windows have $(d_r,d_m,R)=(0.4,0,0.4)$ and $(0,0.4,0.2)$.
They accept $(T_r,T_m)=(0.4,0.2)$. Allocating once to the integrated demands
$(0.4,0.4)$ and receiving $0.6$ gives $(0.3,0.3)$ instead. The latter assigns
early unused mainline opportunity to mainline traffic that arrives later.
This is a conditional node diagnostic, not a reached control witness. A native
frozen-request CTM step is a separately declared numerical transcription.

For each section/lane $b$, let $T_{a b}$ be completed transfers, $B_b$ admitted
external entry and $E_b$ boundary exit. Conditional admissible ledgers satisfy
\begin{equation}
 n_b^{k+1}=n_b^k+B_b+\sum_a T_{a b}-E_b-\sum_c T_{b c},\qquad
 \sum_b(n_b^{k+1}-n_b^k)=\sum_b(B_b-E_b).
\end{equation}
Internal transfers cancel. The same balances apply to the finite approach
and outside waiting. Nonnegativity/storage require compatible sending and
receiving bounds; cancellation alone does not prove a valid traffic solver.

\newpage
\section*{3. Local count coordinates avoid counting an entry twice}
For a section $[0,L]$, all declared lane/ramp transfers enter at its upstream
port. Let $C_b(t)$ be cumulative \emph{all} accepted incoming traffic plus an
initial count gauge. With no interior source in this section, define
\begin{equation}
 N_b(t,x)=C_b(t)-\int_0^x\rho_b(t,\xi)\,d\xi,\qquad
 \partial_xN_b=-\rho_b,\quad \partial_tN_b=F_b(\rho_b).
\end{equation}
The time identity is for the own conservative field, since
$C_b'=F_b(\rho_b(t,0+))$ is its accepted boundary flux. Along a restriction,
\begin{equation}
 \frac{d}{dt}N_b(t,x_j(t))=F_b(\rho_b)-v_j\rho_b=0
\end{equation}
on a compatible zero-relative-passing path. Such a path has a constant local
count anchor until crossing a section boundary or a separately declared source.
At an actual crossing time $t_*$, the new local label is
\begin{equation}
 h_j^{\mathrm{new}}=N_{b'}(t_*,0).
\end{equation}
This rebase changes coordinates, not restriction inventory or traffic mass.
Port entries already occur in $C_{b'}$ and must not also be inserted as an
interior source. A global-lane formulation needs its corresponding source-prefix
corrections. A stationary restriction at a node uses the common receiving
constraint for every incoming stream. These identities do not settle all
physical node/path interactions.

\section*{4. A conservative cell projection can contradict exact records}
Use triangular KW
\begin{equation}
 F(\rho)=\min\{u\rho,w(\kappa-\rho)\},\quad
 u=30,\quad w=5,\quad\kappa=0.15,\quad L=30,\quad\Delta t=1.
\end{equation}
Units are m, s and vehicle equivalents. Initial density is $0.01$ throughout
the cell, with a free outlet and sufficient upstream demand. Two old records
start at $x=h=0$, with desired speeds 5 and 15 m/s. This obeys $u\Delta t=L$.
The established local coincident fan has the slow record holding back traffic,
while the faster record separates through vacuum. Behind the slow record,
\begin{equation}
 \rho_U(5)=\frac{w\kappa}{w+5}=0.075,\qquad q_U(5)=5\rho_U(5)=0.375.
\end{equation}
At $t=1$, $x_1=5$, $x_2=15$ and the original free traffic has exited. The
remaining field is $0.075$ on $[0,5]$ and zero on $(5,30)$. Initial mass is
0.3, inflow is 0.375 and outflow is 0.3. Hence
\begin{equation}
 n(1)=0.3+0.375-0.3=0.375,\quad C(1)=0.375,\quad
 N(1,5)=N(1,15)=N(1,30)=0.
\end{equation}
The initial free-traffic tail moves at 30 m/s, so no conflict with the fast
record is introduced. These are local initial data, not a reached E1/meter
case. Neither record adds traffic population.

\newpage
Uniform cell projection instead gives
\begin{equation}
 \bar\rho=\frac{0.375}{30}=0.0125,\quad
 \bar N(5)=0.3125,\quad\bar N(15)=0.1875,\quad
 \int_5^{15}\bar\rho\,dx=0.125.
\end{equation}
Cell total and boundary counts are correct, but the two old zero-count
anchors disagree with the field. Relabeling them changes source-free count
histories; repositioning them changes selected motion. The conclusion is
about this projection's joint compatibility, not every CTM scheme and not
the impossibility of a qualified numerical error contract.
Indeed for one projection, both densities lie in $[0,\kappa]$, so
$|N(x)-\bar N(x)|\le\kappa L$ follows by integrating their difference.
This bound shrinks as cell length shrinks for a fixed field, but does not
control one-sided traces, successive endogenous updates or derivatives.

\section*{5. Existing anchors permit an exact local reconstruction}
Include cell boundaries and retained record locations, all expressed in the
completed own local count gauge. Sort unique locations $z_0<\cdots<z_m$ with
counts $a_0,\ldots,a_m$. Coincident records must have equal counts; merging
their \emph{field location} does not merge their distinct kinetic metadata.
Compatibility requires
\begin{equation}
 0\le a_i-a_{i+1}\le\kappa(z_{i+1}-z_i)\quad(0\le i<m).
\end{equation}
If it holds, define the candidate F.A field on each interval by
\begin{equation}
 \rho_i=\frac{a_i-a_{i+1}}{z_{i+1}-z_i},\qquad
 N(x)=a_i-\rho_i(x-z_i),\quad z_i\le x\le z_{i+1}.
\end{equation}
Then $0\le\rho_i\le\kappa$, every anchor is matched, and
\begin{equation}
 \int_{z_0}^{z_m}\rho\,dx=\sum_{i=0}^{m-1}(a_i-a_{i+1})=a_0-a_m.
\end{equation}
Thus completed cell mass is unchanged. Equal adjacent counts produce vacuum.
The example's nodes $(0,0.375),(5,0),(15,0),(30,0)$ reproduce its field. With
only cell boundaries, this is ordinary uniform cell reconstruction.

This proves reconstruction only. Invalid inputs require diagnosis, without
clipping or relabeling. Averaging between anchors is still numerical. F.A is
unadopted; F.B needs joint field/count/motion error qualification. K.A stays
fixed. These intervals are not new CTM cells using the old timestep: next
fluxes need their own CFL/stability or appropriate interior wave resolution.

\section*{6. A clearing tail needs more than asymptotic stability}
Finite descriptor births do not imply traffic clearance. Conditional on an
ordinary free-flow single-cell regime with no external entry and positive
progress $\gamma=u\Delta t/L\in(0,1]$, its remaining mass satisfies
\begin{equation}
 n_{k+1}=(1-\gamma)n_k,\qquad
 \Delta t\sum_{k=0}^{\infty}n_k=\frac{\Delta t\,n_0}{\gamma}.
\end{equation}
This is discrete left-endpoint TST, which can differ from the continuous-time
integral. For $\gamma<1$, traffic approaches zero without finite-step emptiness,
yet TST converges. In contrast, $n(t)=1/(1+t)$ tends to zero with divergent
integral. No result proves entry into this free regime. Ordinary CTM theorems
do not certify the coupled task; complete TST, metering response and gradients
remain pending.
\end{document}
